\section*{Exercises}
\markright{\MakeUppercase{Exercises}}%
\subsection*{Exercise 1}

Write a Ymir program that prints the following to the console:
\begin{verbatim}
Hello!
I wrote my first Ymir program.
Ain't that fun?
\end{verbatim}

\subsection*{Exercise 2}
The following program contains several errors:

\begin{lstlisting}[style=coloredverbatimError, escapechar=@, caption=Ymir program with errors]
*/ Comments of the main function @/*@
fn main {
  println("If this text")
  println("appears on the screen, ');
  println('you got it right').
)
\end{lstlisting}
Correct the errors and compile the program with \token{gyc} to test your
corrections.

\subsection*{Exercise 3}

What does the following program output to the console?
\begin{lstlisting}[style=coloredverbatim]
use std::io;

fn pause() {
  print("BREAK");
}

fn main() {
  println("\nDear reader, ");
  print("have a ");
  pause();
  println("!");
}
\end{lstlisting}

\subsection*{Exercise 4}

The following program is correct, and prints \textit{1 + 2 = 3} when it is
executed, but it does not follow the layout presented in
Section~\ref{sec:basics:layout}. Rewrite it so that it does.

%% restyle: skip -- the point of the exercise is that this layout is wrong
\begin{lstlisting}[style=coloredverbatim]
use std::io;
fn add (a : i32,b : i32)-> i32 {a+b}
fn main
(){
let x=add (1 , 2);
println ("1 + 2 = ",x);
}
\end{lstlisting}

\subsection*{Exercise 5}

The compiler was invoked on a source file named \textit{calc.yr}, and produced
the following representation. In which file was this representation written, and
which command line produced it? What does the program print when it is executed?

\begin{lstlisting}[style=myilVerb]
frame : calc::twice (a(#2) : i32)-> i32 return (a(#2) * 2)
frame : calc::main ()-> void {
    std::io::println!{[c8], i32}::println ("= "s8, calc::twice (21));
    <unit-value>
}
\end{lstlisting}

\subsection*{Exercise 6}

Using Gyllir, create a project named \textit{greeter} that prints
\textit{Hello reader!} when it is run. Which commands must be typed, and which
file of the project must be modified?

\vfill%
\pagebreak

\forceevenpage{}

\section*{Solutions}
\markright{\MakeUppercase{Solutions}}%
\subsection*{Exercise 1}

\begin{lstlisting}[style=coloredverbatim, caption=Solution for exercise 1]
use std::io;

fn main() {
  println("Hello!");
  println("I wrote my first Ymir program.");
  println("Ain't that fun?");
}
\end{lstlisting}

\subsection*{Exercise 2}

%% check: skip -- answer listing is written almost entirely in LaTeX escapes
\begin{lstlisting}[style=coloredverbatim, escapechar=@]
@\hcb{\color{purple}{use} \color{black}{std::io;}}@

@\hcb{/*}@ @\textit{Comments of the main function}@ @\hcb{*/}@
fn main@\hcb{()}@ {
  println(@\color{black!30!applegreen}{"If this text"}@)@\hcb{;}@
  println(@\color{black!30!applegreen}{"appears on the screen,}@ @\hcb{"}@);
  println(@\hcb{"}@@\color{black!30!applegreen}{you got it right}@@\hcb{"}@);
@\hcb{\}}@
\end{lstlisting}

\subsection*{Exercise 3}
The output starts with a line break.
\begin{lstlisting}[style=bashVerb]

Dear reader,
have a BREAK!
\end{lstlisting}

\subsection*{Exercise 4}

Each statement is written on its own line, the body of the two functions is
indented by two spaces, the parameter lists and the type colons are attached to
the name that precedes them, and the binary operators are surrounded by spaces.

\begin{lstlisting}[style=coloredverbatim, caption=Solution for exercise 4]
use std::io;

fn add(a: i32, b: i32)-> i32 {
  a + b
}

fn main() {
  let x = add(1, 2);
  println("1 + 2 = ", x);
}
\end{lstlisting}

\subsection*{Exercise 5}

This is a M-YIL representation, so it was written in the file
\textit{calc.yr.ydump-sem}, by the command line below. The L-YIL
representation of the same source file was written in \textit{calc.yr.ydump-yil}
at the same time.

\begin{lstlisting}[style=bashVerb, escapechar=@+]
@+\textcolor{teal!80}{alice@dev:\mtilde}@+$ gyc -fdump-ymir calc.yr -o calc
\end{lstlisting}

The \token{main} function calls \token{println} with two values, the string
\token{"= "} and the result of the call to \token{twice} with the value
\token{21}. The frame of \token{twice} returns its parameter multiplied by two,
so the program prints \textit{= 42}.

\subsection*{Exercise 6}

The \textit{init} subcommand creates the project, whose name is asked by the
first question. The file to modify is \textit{src/main.yr}, the source file
created by Gyllir, which prints \textit{Hello World!} until it is changed.

\begin{lstlisting}[style=bashVerb, escapechar=@+]
@+\textcolor{teal!80}{alice@dev:\mtilde}@+$ mkdir greeter
@+\textcolor{teal!80}{alice@dev:\mtilde}@+$ cd greeter
@+\textcolor{teal!80}{alice@dev:\mtilde}@+@+\textcolor{green!60!black}{/greeter}@+$ gyllir init
Name [main]: greeter
@+\textcolor{black!50}{... the other questions can keep their default value}@+
\end{lstlisting}

\begin{lstlisting}[style=coloredverbatim, caption=Content of \textit{src/main.yr}]
use std::io;

fn main() {
  println("Hello reader!");
}
\end{lstlisting}

The project is then compiled and executed by the \textit{run} subcommand.

\begin{lstlisting}[style=bashVerb, escapechar=@+]
@+\textcolor{teal!80}{alice@dev:\mtilde}@+@+\textcolor{green!60!black}{/greeter}@+$ gyllir run
  @+\color{black!30!applegreen}{Compiling}@+ greeter v0.1.0
  @+\color{black!30!applegreen}{Compiling}@+ 1 modules in thread 1 (among 1 changed files of 1)
  @+\color{black!30!applegreen}{Produced}@+ executable file greeter
  @+\color{black!30!applegreen}{Executing}@+ project greeter v0.1.0

Hello reader!
\end{lstlisting}
